\documentclass[final,3p,times,11pt,authoryear]{elsarticle}

% Preambulo, autores, afiliacoes e declaracoes: template do EJOR fornecido pelos autores (Novo_LCPP_EJOR).
% Este arquivo e o MODELO usado por make_latex.py; main.tex e GERADO (nao edite main.tex: edite o modelo ou o manuscrito).

\usepackage[utf8]{inputenc}
\usepackage[T1]{fontenc}
\usepackage{amsmath,amssymb,amsfonts,amsthm}
\usepackage{graphicx}
\usepackage{textcomp}
\usepackage{subcaption}
\usepackage{booktabs}
\usepackage{longtable}
\usepackage{multirow}
\usepackage{comment}
\usepackage{url}
\usepackage{float}
\usepackage{algorithm}
\usepackage{algorithmic}
\floatstyle{ruled}
\restylefloat{algorithm}
% \usepackage{cuted}  % do template; nao usado aqui e ausente desta instalacao do MiKTeX
\usepackage{tabularx}
\usepackage{setspace}
\onehalfspacing
\usepackage{enumitem}
\setlist{nosep,topsep=2pt,itemsep=1pt,parsep=0pt,leftmargin=*}
\setlength{\textfloatsep}{8pt plus 2pt minus 2pt}
\setlength{\floatsep}{6pt plus 2pt minus 2pt}
\setlength{\intextsep}{6pt plus 2pt minus 2pt}
\setlength{\abovecaptionskip}{4pt}
\setlength{\belowcaptionskip}{2pt}
\usepackage{xcolor}
\usepackage{bm}
\usepackage{hyperref}
\hypersetup{colorlinks=true,linkcolor=blue!50!black,citecolor=blue!50!black,urlcolor=blue!50!black}
\emergencystretch=3em

\usepackage{seqsplit}

\newcommand{\code}[1]{\texttt{\seqsplit{#1}}}

%% Theorem-like environments
\theoremstyle{plain}
\newtheorem{proposition}{Proposition}
\newtheorem{lemma}{Lemma}
\newtheorem{theorem}{Theorem}
\newtheorem{corollary}{Corollary}
\theoremstyle{definition}
\newtheorem{definition}{Definition}
\newtheorem{assumption}{Assumption}
\theoremstyle{remark}
\newtheorem{remark}{Remark}

\newcommand{\todo}[1]{\textcolor{red}{\textbf{[TODO #1]}}}
\newcommand{\pending}[1]{\textcolor{orange}{\textit{[PENDING #1]}}}

\journal{European Journal of Operational Research}

\begin{document}
\begin{frontmatter}

\title{%%TITLE%%}

%% Authors
\author[uece]{Bonfim Amaro Junior\corref{cor1}}
\ead{bonfim.amaro@uece.br}
\cortext[cor1]{Corresponding author}
\author[ufmg]{Marcio Costa Santos}
\ead{marciocs@dcc.ufmg.br}
\author[uece]{Guilherme Nepomuceno de Carvalho}
\ead{nepomuceno.carvalho@aluno.uece.br}
\author[unifor]{Plácido Rogério Pinheiro}
\ead{placido@unifor.br}
\author[uece]{Carlos Vinicio Dantas Araújo}
\ead{carlos.vinicio@aluno.uece.br}

%% Affiliations
\affiliation[uece]{organization={Universidade Estadual do Ceará (UECE)},
            addressline={Av. Dr. Silas Munguba, 1700},
            city={Fortaleza},
            postcode={60714-903},
            state={Ceará},
            country={Brazil}}
\affiliation[ufmg]{organization={Universidade Federal de Minas Gerais (UFMG), Departamento de Ciência da Computação},
            addressline={Av. Pres. Antônio Carlos, 6627},
            city={Belo Horizonte},
            postcode={31270-901},
            state={Minas Gerais},
            country={Brazil}}
\affiliation[unifor]{organization={Universidade de Fortaleza (UNIFOR)},
            addressline={Av. Washington Soares, 1321},
            city={Fortaleza},
            postcode={60811-341},
            state={Ceará},
            country={Brazil}}

\begin{highlights}
%%HIGHLIGHTS%%
\end{highlights}

\begin{abstract}
%%ABSTRACT%%
\end{abstract}

\begin{keyword}
%%KEYWORDS%%
\end{keyword}

\end{frontmatter}

%%SECTIONS%%

%%
%% Sections required by EJOR
%%
\begingroup
\small
\begin{spacing}{1.0}

\section*{CRediT authorship contribution statement}
% Papeis copiados do template (artigo de corte a laser): os autores devem conferir e ajustar a este trabalho.
\textbf{Bonfim Amaro Junior} (Conceptualization, Methodology, Software, Writing original draft).
\textbf{Marcio Costa Santos} (Methodology, Formal analysis, Writing review and editing).
\textbf{Guilherme Nepomuceno de Carvalho} (Software, Validation, Investigation).
\textbf{Carlos Vinicio Dantas Araújo} (Software, Validation, Investigation, Data curation, Writing review and editing).
\textbf{Plácido Rogério Pinheiro} (Supervision, Writing review and editing, Funding acquisition).

\section*{Declaration of Competing Interest}
The authors declare that they have no known competing financial interests or personal relationships that could have appeared to influence the work reported in this paper.

\section*{Declaration of Generative AI and AI-Assisted Technologies in the Writing Process}
%%AIDECL%%

\section*{Data and Code Availability}
All instances are public: the MDPLIB~2.0 library \citep{marti2021a}, the CDP benchmark of \citet{peiro2021} and the GDP benchmark of \citet{lozanoosorio2022} (Zenodo record 18233882, CC-BY 4.0). The source code, the run files (one JSON file per instance, method and replicate, with seed, configuration hash and Pareto front), the frozen reference fronts and the scripts that regenerate every table and figure are available in the public repository \url{https://github.com/CarlosAraujo7/moead-aos-bdp} (code under the MIT licence; data and results under CC-BY 4.0). The file \texttt{README\_REPLICATION.md} of the repository lists the commands that regenerate each table and figure, and \texttt{regression\_check.py} re-executes stored runs from their seeds and checks bit-for-bit equality. Third-party material (the benchmark instances and the published solutions used for the comparison in Section~\ref{sec:anchor_lit}) is credited in the repository README. \todo{archive a release of the repository on Zenodo and add its DOI here}

\section*{Funding}
This work was supported by FAPEMIG-APQ-03509-26, CNPq 448503/2025-5, CNPq 442977/2023-9, CNPq 444528/2024-5, CNPq 421005/2025-4 and by CNPq, PRPq, UFMG via PIBIC fellowship.

\section*{Acknowledgements}
The authors thank the institutions Universidade Estadual do Ceará (UECE), Universidade Federal de Minas Gerais (UFMG), and Universidade de Fortaleza (UNIFOR) for institutional support during the
development of this work.

\end{spacing}
\endgroup

%%
%% References (author-date style via natbib/elsarticle-harv)
%%
\begingroup
\scriptsize
\setlength{\bibsep}{0pt}
\begin{spacing}{0.87}
\bibliographystyle{elsarticle-harv}
\bibliography{referencias}
\end{spacing}
\endgroup

\end{document}
